\section{Spiritual Beings}

In esoteric cosmologies, in the transition from essence to existence, the being passes through a series of influences, from the subtle to the dense. According to \textbf{John Woodroffe}, in Tantric teachings the schema of the Tattvas from the Samkhya school is important in explaining the various states of the being. In the West, Hermetism has an analagous teaching. John Fludd describes the cosmic process\footnote{\url{https://www.meditationsonthetarot.com/the-cosmic-hierarchy}} as a series of 22 states from the Divine Mind to Earth. We have adapted his diagram\footnote{Special thanks to Lydia Bisanti (\url{https://gornahoor.net/?p=1497}) for creating the diagram.} to illustrate this clearly. First of all, we point out that these states can be categorized in three Realms corresponding to the Traditional division of Heaven, Man, and Nature. Each of these realms relate to states of beings and to a particular Hermetic science.

\begin{table}[h]\small\centering
\begin{tabular}{ccc}\toprule
\textbf{Realm} & \textbf{States of Being} & \textbf{Hermetic Science} \\\midrule
Heaven (2 to 10) & Transhuman states & Theurgy or Sacred Magic \\
Man (11 to 18) & Psychological states & Astrology \\
Nature (19 to 22) & Subconscious states & Alchemy \\\bottomrule
\end{tabular}
\end{table}

Note how is integrating the pagan deities into a larger continuum. This is another example of the unity of the Roman Tradition. There are a few main points about this diagram.

\begin{description}

\item[Unity at the top, Diversity at the bottom]

While there is an ontologically unity in principle, at the level of manifestation there is wide diversity. All attempts to establish a unity at the level of manifestation have failed and will continue to fail. But the modern flatland mind can only understand unity at that level, and will extrapolate backwards in time to a fictitious source at the beginning.
\item[Inner experience]

It is through inner experience that we come to know and understand these forces. They are not proved by philosophy nor by science.
\item[Represent states of being]

The layers correspond to states of the being, subhuman, human and transhuman. We may sense the obstacles formed by elemental forces, or the forces unleashed by the gods as alien to us. But as Evola points out, those experiences are actually a privation, that is, a lack in us, not independent of us. As we gain power, we begin to understand them as our own state.
\item[Personal not impersonal : the I]

It is a sign of ``sophistication'' today to regard the gods or God as impersonal forces. However, the Absolute is the I, which acts by its Will, just the opposite of an impersonal force.
\item[Do not act directly]

These states do not act directly on the world, but rather with man as intermediary\footnote{\url{https://gornahoor.net/?p=1798}}. Each layer contributes through thoughts, ideas, and tendencies.
\end{description}


\paragraph{Nature}


We have already addressed this issue here: Possession by Evil Spirits\footnote{\url{https://www.meditationsonthetarot.com/possession-by-evil-spirits}}, so there is no need to repeat it.

\paragraph{Man}


In his sermon on Mars Hill in Athens, Paul criticized the belief in the power of idols. What this really means is the acknowledgment of the end of paganism as an exoteric practice. Nevertheless, Paul does quote some pagan poets, which is \emph{ipso facto} an admission of their inspiration. This was commonly accepted in Hellenic Christianity until the doctrine of \emph{sola scriptura} undermined any possibility of that in the West. So if Homer is inspired, then the myths of the gods necessarily refer to an esoteric teaching. Hence, the gods are symbolic figures and we recognize their influence inner experience.

More than a millennium later, Robert Fludd\footnote{\url{https://en.wikipedia.org/wiki/Robert_Fludd}} had the same understanding. He recognized the power of the gods and incorporated them into the Chain of Being. But it was only in the 20th century that \textbf{Carl Jung} began to explore the influences of these forces in a systematic way and to bring their activities from the unconscious into consciousness.

\paragraph{Heaven}


Beyond the psychological states in man, there are higher transformal states. These have been called angels in several traditions and this particular schema is due to Dionysius who was a Greek associate of Paul. That relationship made his teaching on the angels authoritative in the Middle Ages. A more recent work on the angels is \emph{Spiritual Beings} by \textbf{Rudolf Steiner}. Rene Guenon recognizes the existence of the celestial hierarchy, but doesn't dwell on the topic since, from a metaphysical point or view, it is sufficient to note the existence of higher states, but not the specifics. However, for the Hermetist, whose project is to achieve such higher states, knowledge of this hierarchy is important.

In passing, we can tie this back to Julian. The archangels, in particular, are the forces behind the different nations. Steiner writes: ``A nation is a homogeneous group of people directed by one of the Archangels.'' This used to be the common belief in Europe. The different ranks of angels ``guide'', in the sense of forming ideas and worldviews, various aspects of life. That is why similar constellations of related ideas can be held by a group of people almost in unison. The existence of a counter-hierarchy is left for another time.


\flrightit{Posted on 2011-03-17 by Cologero}

